\section{Experimental Results}
\label{sec:results}
All numbers in this section are read from generated tables and figures; the result tables are written by the statistics scripts and are not
retyped. Of the Study~I hypotheses, H1$'$ (dose-response) passes; H2$'$ passes for bone-region PSNR against the reference model and against
the no-twin control and fails for SSIM, as the analysis plan anticipated; H3$'$ passes for both contrasts on VinDr-RibCXR; H4$'$a is mixed
(ours above the reference model on JSRT and below it on Node21, above the no-twin control on both); H4$'$b shows no downstream benefit of
the soft-tissue image used alone; and H5$'$ holds on synthetic data with exact truth and fails under the shift to real radiographs. The v1
hypothesis H1 of the first analysis plan passes, and in the exploratory site calibration (SAP\_v5) H16--H18 pass. In Study~II the primary
hypothesis H6 fails on all three datasets, and in Study~III the primary hypothesis H11 fails. All analysis changes made after the freezes are listed in the deviations log of the repository.

% ---------------------------------------------------------------------------------------------------------------------
\subsection{Paired Dual-Energy Fidelity, Dose-Response and Ablation (Study~I: H1\texorpdfstring{$'$}{'}, H2\texorpdfstring{$'$}{'})}
\label{sec:res-fidelity}

\subsubsection{Fidelity on held-out paired radiographs (H2$'$)}
On the 51 held-out SZCH-X-Rays pairs (Table~\ref{tab:szch_heldout}) the bone-region PSNR$_{\mathrm{cal}}$ of ours is 21.46\,dB and that of
the reference model 21.51\,dB. The paired difference is $-0.06$\,dB with a one-sided 95\,\% lower bound of $-0.15$\,dB, inside the
non-inferiority margin of $-0.5$\,dB (Table~\ref{tab:h2}). Against the no-twin control the difference is $+0.01$\,dB (lower bound
$-0.02$\,dB, margin $-0.25$\,dB): the twin regulariser costs no measurable bone-region fidelity. SSIM$_{\mathrm{cal}}$ fails its margin: ours
is lower by 0.021 (lower bound $-0.023$, margin $-0.01$). The gap sits in fine texture. With a band limit of $\sigma=4$ the SSIM of ours and
of the reference model are 0.946 and 0.948, and the high-frequency energy of the non-bone region, relative to the dual-energy soft-tissue
image, is 1.09 for ours and 0.99 for the reference model; LPIPS is 0.234 and 0.166. BS-LDM, which was trained on dual-energy pairs from the
same hospital, reaches 28.57\,dB; it is shown for orientation and is not part of a pre-specified comparison. Fig.~\ref{fig:pareto} places
fidelity against lesion absorption, and Fig.~\ref{fig:qualitative} shows three held-out cases.

\begin{table*}[!t]
\centering
\caption{Held-out SZCH-X-Rays pairs: bone-region and whole-image metrics for ours and the comparators, calibrated (primary) and raw. HF
ratio: high-frequency energy of the non-bone region relative to the dual-energy soft-tissue image (1 = equal).}
\label{tab:szch_heldout}
\restable[\textwidth]{szch_heldout}
\end{table*}

\begin{table*}[!t]
\centering
\caption{H2$'$ non-inferiority on the 51 held-out SZCH-X-Rays pairs (paired differences, image bootstrap, one-sided 95\,\% lower bound).}
\label{tab:h2}
\restable[\textwidth]{h2_noninferiority}
\end{table*}

\begin{figure}[!t]
\centering
\safefig[\columnwidth]{figures/pareto.pdf}
\caption{Lesion absorption on VinDr-RibCXR against bone-region PSNR$_{\mathrm{cal}}$ on the 51 held-out SZCH-X-Rays pairs. Line: the
SwinV2-T family at twin weights $w=0\ldots64$ (single seeds light, seed means dark). Markers: ours and the no-twin control
(ConvNeXt-V2-T, $K=3$), the reference model, the v1 model and the identity. Lower absorption at equal fidelity is better.}
\label{fig:pareto}
\end{figure}

\begin{figure*}[!t]
\centering
\safefig[0.8\textwidth]{figures/qualitative_szch.png}
\caption{Held-out SZCH-X-Rays cases (CC BY-NC 4.0 \cite{sun2026bsldm}): radiograph, dual-energy soft-tissue reference, the soft-tissue
estimates of ours and of the reference model, and the bone estimate of ours. The three cases are not hand-picked: they are the held-out
images at the 25th, 50th and 75th percentile of the bone-region PSNR$_{\mathrm{cal}}$ of ours.}
\label{fig:qualitative}
\end{figure*}

\subsubsection{Site-calibrated fidelity (exploratory, H16--H18)}
\label{sec:res-sitecal}
The smooth site correction of Section~\ref{sec:m-sitecal}, fitted on the 24 development pairs, raised the bone-region PSNR$_{\mathrm{cal}}$
of ours on the 51 held-out pairs from 21.45 to 23.58\,dB, a paired gain of 2.12\,dB (one-sided 95\,\% lower bound 1.99\,dB; H16 passes),
and its whole-image PSNR$_{\mathrm{cal}}$ from 25.56 to 28.30\,dB (gain 2.74\,dB, lower bound 2.61\,dB; H16b passes;
Tables~\ref{tab:sitecal_fid} and~\ref{tab:sitecal_tests}). Every CT-supervised model gained about as much (reference model 21.51 to
23.92\,dB, v1 model 21.95 to 24.06\,dB, no-twin control 21.45 to 23.74\,dB) and the unprocessed radiograph less (20.21 to 21.46\,dB): the
correction calibrates the site rather than improving a suppressor, and the order of the CT-supervised models did not change. After
calibration ours remained non-inferior to the calibrated reference model, by a narrower margin than before ($-0.35$\,dB, lower bound
$-0.44$\,dB against the margin of $-0.5$\,dB; H17 passes), and the no-twin control was slightly ahead of ours (23.74 against 23.58\,dB; not
tested).

The calibration did not cost lesion fidelity (Table~\ref{tab:sitecal_ribcxr}): on VinDr-RibCXR the calibrated ours absorbed 0.036 of an
inserted lesion and the calibrated reference model 0.132 (difference $-0.096$ $[-0.105,-0.087]$; H18 passes), and the calibration lowered
the absorption of both models slightly ($-0.015$ and $-0.017$); the uncalibrated models of this run reproduced the audit of
Table~\ref{tab:audit_ribcxr}. It weakened suppression a little and made it more selective: the rib-edge reduction of ours fell from 0.298 to
0.282 and the reduction on other lung pixels from 0.078 to 0.050, so the selectivity rose from 0.220 to 0.232. BS-LDM, trained on this
site's dual-energy pairs, remains about 5\,dB ahead in the bone region (28.57 against 23.58\,dB) and 3.65\,dB in the whole image.

\begin{table*}[!t]
\centering
\caption{Site-calibrated fidelity on the 51 held-out SZCH-X-Rays pairs (exploratory, SAP\_v5): every method as is and with the smooth site
correction of Eq.~\eqref{eq:sitecal} fitted on the 24 development pairs. BS-LDM, trained on this site's dual-energy pairs, is not
corrected.}
\label{tab:sitecal_fid}
\restable[\textwidth]{site_cal_fidelity}
\end{table*}

\begin{table*}[!t]
\centering
\caption{SAP\_v5 hypotheses (exploratory): paired differences with image bootstrap (H16, H16b, H17: one-sided 95\,\% lower bounds) or
image-cluster bootstrap intervals (absorption on 1,470 VinDr-RibCXR sites), \num{10000} resamples.}
\label{tab:sitecal_tests}
\restable[\textwidth]{site_cal_tests}
\end{table*}

\begin{table*}[!t]
\centering
\caption{VinDr-RibCXR (metrics only): mean lesion absorption of one audit run with all four models (shared sites) and the rib benchmark,
as is and site-calibrated (exploratory, SAP\_v5).}
\label{tab:sitecal_ribcxr}
\restable[\textwidth]{site_cal_ribcxr}
\end{table*}

\subsubsection{Dose-response of the twin regulariser (H1$'$)}
In the SwinV2-T family (3, 2, 2, 4, 2 and 2 seeds at $w=0,4,8,16,32,64$; Table~\ref{tab:dose_response}, Fig.~\ref{fig:dose_response}) the
mean absorption on VinDr-RibCXR falls from 0.183 at $w=0$ to 0.043 at $w=64$. The slope of $A$ on $\log_2(1+w)$ is $-0.0233$
$[-0.0249,-0.0217]$; its upper bound is below zero, so H1$'$ passes. On the Qure-projection test split with exact truth the slope is
$-0.0174$ $[-0.0190,-0.0160]$, with $A$ levelling off above $w=16$ (0.118, 0.122 and 0.112 at $w=16,32,64$). Bone-region
PSNR$_{\mathrm{cal}}$ on the held-out pairs stays between 21.22 and 21.28\,dB up to $w=16$ and drops to 20.92 and 21.01\,dB at $w=32$ and
64, which is why $w=16$ is the working point (Fig.~\ref{fig:pareto}).

\begin{table}[!t]
\centering
\caption{Dose-response of the lesion absorption on the twin weight (SwinV2-T family): seeds per weight, mean $A$ on VinDr-RibCXR and on the
Qure-projection test split, and the pre-specified slope with its image-cluster bootstrap interval. Each weight averages all of its seeds
(four at $w=16$, seeds 0--3); the encoder comparisons use seeds 0--2 only.}
\label{tab:dose_response}
\restable{dose_response}
\end{table}

\begin{figure}[!t]
\centering
\safefig[\columnwidth]{figures/dose_response.pdf}
\caption{Dose-response of the lesion absorption $A$ on VinDr-RibCXR on the twin weight $w$ (SwinV2-T), one marker per seed, with the
pre-specified slope on $\log_2(1+w)$.}
\label{fig:dose_response}
\end{figure}

\subsubsection{Ablation of the regulariser}
The ablation (Table~\ref{tab:ablation_twin}) is development evidence: it uses the audit on all 75 SZCH images and the 24 development pairs.
At $w=16$ the full regulariser gives $A=0.067$ against 0.190 without it. Three variants lose the effect: a stop-gradient on the branch
without the pasted lesion (0.185), a consistency term over the whole image instead of the lesion support (0.197), and lesion insertion
without the twin term (CLIA, 0.191). Dilating the support by 31\,px keeps part of it (0.137); up to four lesions per twin instead of up to
two keep all of it (0.074). The effect therefore needs gradients through both branches and a consistency term confined to where the lesion
was pasted; showing the model lesions during training is not enough.

\begin{table*}[!t]
\centering
\caption{Ablation of the twin regulariser at $w=16$ (SwinV2-T, 300k samples, mean over seeds): one factor changed per row. Development
evidence: audit on all 75 SZCH images, bone PSNR$_{\mathrm{cal}}$ on the 24 development pairs.}
\label{tab:ablation_twin}
\restable[\textwidth]{ablation_twin}
\end{table*}

\subsubsection{The v1 hypothesis H1}
The v1 hypothesis H1 of the first plan was tested on synthetic LIDC test CTs with exact truth on 27 October 2026, three seeds per arm
(Supplementary Table~S9), and passes: CLIA absorbs less than the recipe arm ($-0.022$ $[-0.032,-0.013]$, Holm $p<0.001$) with non-inferior
bone MAE. The lesion-weighted loss also lowers absorption ($-0.019$) and is non-inferior on the test CTs, although it missed the margin on
the validation CTs that decided the selection; combining both lowers absorption most ($-0.066$) but misses the margin. The 300k-sample finals
of CLIA and recipe replicate the contrast ($-0.024$ $[-0.034,-0.014]$). On real radiographs the v1 models remain far above the v2 models
($A=0.289$ and 0.308 on VinDr-RibCXR).

% ---------------------------------------------------------------------------------------------------------------------
\subsection{Confirmatory Counterfactual Audit, Rib Benchmark and Real-Nodule Contrast (Study~I: H3\texorpdfstring{$'$}{'},
H4\texorpdfstring{$'$}{'}a)}
\label{sec:res-audit}

\subsubsection{Confirmatory audit on VinDr-RibCXR (H3$'$)}
On 1,470 sites in 245 VinDr-RibCXR images (Tables~\ref{tab:audit_ribcxr} and~\ref{tab:audit_contrasts}) ours absorbs $A=0.051$
$[0.045,0.057]$ of an inserted lesion, the reference model 0.149 $[0.138,0.160]$ and the no-twin control 0.178 $[0.167,0.188]$. The paired
differences are $-0.098$ $[-0.108,-0.089]$ against the reference model and $-0.127$ $[-0.134,-0.119]$ against the no-twin control
(Holm-adjusted $p<0.001$ for both), so H3$'$ passes. The advantage does not depend on rib overlap here: the interaction between sites with at
least 15\,\% overlap with an annotated rib and the other sites is $0.003$ $[-0.017,0.027]$ against the reference model. On JSRT (secondary)
the pattern repeats ($A$ 0.048, 0.118 and 0.121; differences $-0.070$ and $-0.073$) with a larger advantage at bone-overlapping sites
(interaction $-0.060$ $[-0.074,-0.046]$); there the overlap was defined by the reference model's bone map, not by the consensus map the plan
names. The SZCH audit (development evidence) agrees ($A$ 0.054, 0.130 and 0.195). Fig.~\ref{fig:audit_maps} shows one site.

\begin{table}[!t]
\centering
\caption{Counterfactual lesion audit on VinDr-RibCXR: mean absorption per method with image-cluster bootstrap intervals (lower is
better).}
\label{tab:audit_ribcxr}
\restable{audit_ribcxr}
\end{table}

\begin{table*}[!t]
\centering
\caption{H3$'$ paired site-level contrasts on the three audit sets, Holm over the two contrasts, and the interaction with bone overlap
($\geq$15\,\% vs.\ $<$15\,\%). SZCH rows: the 51 held-out images (306 sites), still development evidence because the audit on all 75 images
was used during development; the JSRT overlap comes from the reference model's bone map.}
\label{tab:audit_contrasts}
\restable[\textwidth]{audit_contrasts}
\end{table*}

\begin{figure*}[!t]
\centering
\safefig[\textwidth]{figures/audit_maps.pdf}
\caption{The audit on one held-out SZCH-X-Rays image (CC BY-NC 4.0): a 7\,mm nodule from a LIDC test CT is added to the radiograph; the
change of each bone estimate (red = lesion signal moved into the bone image) and the resulting soft-tissue image. Not hand-picked: eight
placements over the bone-overlapping lung field were drawn with a fixed seed and the one with the median $A(\text{ours})-A(\text{reference})$
is shown.}
\label{fig:audit_maps}
\end{figure*}

\subsubsection{Rib benchmark}
Absorption must be read beside suppression strength, because doing nothing gives $A=0$ (Table~\ref{tab:rib_benchmark}). The reference model
removes more edge energy, both on the rib boundaries (48\,\% against 30\,\% for ours) and on the other lung pixels (22\,\% against 8\,\%);
its selectivity is 0.264 against 0.220, and the AUC separating rib from non-rib pixels is the same (0.671 and 0.672). Ours suppresses ribs
less strongly, but its absorption is a third of the reference model's while its rib-edge removal is about 60\,\% of it, and its bone-region
fidelity is non-inferior (H2$'$).

\begin{table}[!t]
\centering
\caption{Rib benchmark on VinDr-RibCXR (metrics only): gradient-energy reduction on rib-boundary pixels and on other lung pixels,
selectivity and AUC.}
\label{tab:rib_benchmark}
\restable{rib_benchmark}
\end{table}

\subsubsection{Real nodules without lesion insertion (H4$'$a)}
Real nodules give a mixed answer (Table~\ref{tab:nodule_cnr}). On JSRT (154 nodules) the contrast-to-noise ratio rises by 0.041 with ours
and by 0.018 with the reference model; the difference is $+0.023$ $[0.008,0.041]$, and $+0.032$ $[0.014,0.052]$ for the 123 nodules with at
least 15\,\% bone overlap. On Node21 (1,474 boxes) the reference model raises contrast more: $+0.103$ against $+0.078$, a difference of
$-0.025$ $[-0.033,-0.018]$, and $-0.017$ $[-0.026,-0.009]$ for the 1,109 bone-overlapping boxes. Against the no-twin control ours raises
real-nodule contrast in all four comparisons ($+0.007$ to $+0.027$, all intervals above zero; Holm within data set $p<0.001$).

\begin{table*}[!t]
\centering
\caption{Contrast-to-noise change of real nodules on JSRT and Node21 (soft-tissue image minus original), all lesions and lesions with at
least 15\,\% bone overlap; image-cluster bootstrap intervals.}
\label{tab:nodule_cnr}
\restable[\textwidth]{nodule_cnr}
\end{table*}

\subsubsection{Encoders (exploratory)}
The twin effect replicates across three encoders trained under the same recipe (Table~\ref{tab:encoders}). On VinDr-RibCXR the twin lowers
$A$ by 0.127 for ConvNeXt-V2-T, 0.114 for SwinV2-T and 0.112 for RAD-DINO, and on JSRT by 0.073, 0.061 and 0.064. ConvNeXt-V2-T was chosen
as ``ours'' on development data before the freeze (Section~\ref{sec:m-encoders}); RAD-DINO reaches the same absorption ($A=0.053$ on
VinDr-RibCXR).

\begin{table}[!t]
\centering
\caption{Twin regulariser per encoder (exploratory): absorption difference to the matched no-twin control and its interaction with bone
overlap; three seeds per model (seeds 0--2), SZCH on its 51 held-out images.}
\label{tab:encoders}
\restable{encoders}
\end{table}

% ---------------------------------------------------------------------------------------------------------------------
\subsection{Downstream Detection, Classification and Conformal Uncertainty Transfer (Study~I: H4\texorpdfstring{$'$}{'}b,
H5\texorpdfstring{$'$}{'})}
\label{sec:res-downstream}

\subsubsection{Downstream tasks (H4$'$b)}
The soft-tissue image used alone does not help downstream models (Tables~\ref{tab:downstream_cls} and~\ref{tab:downstream_det},
Fig.~\ref{fig:downstream_forest}). On NIH ChestX-ray14 the macro-AUROC of the three-seed ensemble is 0.818 with the original image and
0.817 with ours ($-0.001$ $[-0.004,0.003]$, patient-cluster bootstrap); per class, ours is lower for effusion, infiltration and pleural
thickening ($-0.004$, $-0.007$, $-0.013$) and higher for nodule ($+0.011$; Holm-adjusted DeLong $p<0.05$ for these four). Adding the soft
image as a second channel helps slightly ($+0.006$ $[0.003,0.008]$), as does the reference model's soft image ($+0.004$ $[0.002,0.007]$). In
detection the change in CPM with ours alone is $+0.001$ $[-0.023,0.023]$ on Node21, $+0.006$ $[-0.008,0.020]$ on TBX11K and $-0.009$
$[-0.016,-0.002]$ on VinBigData; the three-channel input (original, soft and bone) gains $+0.028$ $[0.009,0.047]$ on Node21 and $+0.008$
$[0.002,0.014]$ on VinBigData, and the reference model's soft image gains $+0.029$ $[0.016,0.043]$ on TBX11K. No difference in differences
of the recall at two false positives per image between bone-overlapping and other lesions excludes zero (Node21 $-0.005$
$[-0.071,0.061]$, VinBigData $+0.009$ $[-0.034,0.053]$). On TBX11K it cannot be estimated, because every test box overlaps bone by at least
20\,\%.

\begin{table*}[!t]
\centering
\caption{Downstream classification on NIH ChestX-ray14: macro-AUROC of the three-seed ensemble, per-seed mean $\pm$ SD and the difference
to the original image (patient-cluster bootstrap).}
\label{tab:downstream_cls}
\restable[\textwidth]{downstream_cls}
\end{table*}

\begin{table*}[!t]
\centering
\caption{Downstream detection on Node21, TBX11K and VinBigData: CPM and mAP@50 (mean $\pm$ SD over three seeds), differences to the
original image and the difference in differences (DiD) of recall at two false positives per image (image bootstrap). both: the original,
soft-tissue and bone images as three channels.}
\label{tab:downstream_det}
\restable[\textwidth]{downstream_det}
\end{table*}

\begin{figure}[!t]
\centering
\safefig[\columnwidth]{figures/downstream_forest.pdf}
\caption{Downstream effects of each input against the original image with 95\,\% intervals: macro-AUROC on NIH ChestX-ray14, CPM on the
detection sets.}
\label{fig:downstream_forest}
\end{figure}

\subsubsection{Reliability of the bone-image uncertainty (H5$'$)}
On the Qure-projection test split with exact truth the split-conformal intervals of the three-model ensemble cover 0.900 $[0.897,0.903]$ of
the pixels at a nominal 0.90, and the ensemble spread ranks the absolute bone error well (Spearman 0.76; Table~\ref{tab:reliability},
Fig.~\ref{fig:reliability}). On the 51 held-out SZCH-X-Rays pairs, measured against a pseudo-truth bone image from the dual-energy pair,
the same intervals cover only 0.139 $[0.134,0.145]$ and the spread ranks the error weakly (Spearman 0.11 $[0.07,0.15]$). The uncertainty is
calibrated in the training domain and does not transfer to real radiographs without recalibration.

\begin{table*}[!t]
\centering
\caption{Ensemble spread and split-conformal intervals ($K=3$, nominal coverage 0.90), with image-bootstrap intervals.}
\label{tab:reliability}
\restable[\textwidth]{reliability}
\end{table*}

\begin{figure*}[!t]
\centering
\safefig[0.78\textwidth]{figures/reliability.pdf}
\caption{Left: per-image coverage of the conformal intervals on the Qure-projection test split and on the held-out SZCH-X-Rays pairs (bars:
means; dashed: nominal 0.90). Right: per-image Spearman correlation of ensemble spread and absolute error on SZCH-X-Rays (line: mean).}
\label{fig:reliability}
\end{figure*}

\subsubsection{Recalibration on real pairs (exploratory, post hoc)}
An exploratory analysis, added after this result was known and therefore post hoc on already-seen pairs, recalibrated the intervals on the
24 SZCH-X-Rays development pairs and tested them on the 51 held-out pairs (Table~\ref{tab:bone_recal}, Supplementary Fig.~S10). Pooled
coverage was 0.886 with one global calibration and 0.895 with a separate calibration per stratum of predicted bone overlap and ensemble
spread (Mondrian), whose intervals were less than half as wide (mean width 0.215 against 0.447). A few real pairs thus restore nominal
coverage on average, while the coverage of single images still varies.

\begin{table*}[!t]
\centering
\caption{Exploratory recalibration on real pairs (post hoc): bone-image intervals calibrated on the 24 SZCH-X-Rays development pairs and
tested on the 51 held-out pairs at nominal 0.90. Mondrian: calibration per stratum of predicted bone overlap and ensemble spread.}
\label{tab:bone_recal}
\restable[\textwidth]{v3_bone_intervals}
\end{table*}

% ---------------------------------------------------------------------------------------------------------------------
\subsection{External Validation on Paediatric Radiographs (BDCXR-3257: E1--E3)}
\label{sec:res-bdcxr}
On 3,257 paediatric radiographs from a hospital in Dhaka (Table~\ref{tab:bdcxr_cls}), five-fold cross-validated pneumonia classification
(normal, bacterial, viral) reaches a pooled out-of-fold macro-AUROC of 0.814 on the original image and 0.816 on the soft-tissue image of
ours. The difference, $+0.002$ $[-0.005,0.010]$, meets the pre-specified non-inferiority margin of 0.01 and shows no superiority (E1); the
reference model's soft image gives $+0.001$ $[-0.007,0.010]$ and CLAHE $-0.006$ $[-0.012,0.001]$, which misses the margin. In the lesion
audit on 300 of these images (E2, Table~\ref{tab:bdcxr_audit}) ours absorbs $A=0.084$, the reference model 0.217 and the no-twin control
0.251 (differences $-0.134$ $[-0.144,-0.124]$ and $-0.167$ $[-0.175,-0.160]$), with a larger advantage at bone-overlapping sites. The
inserted lesions come from adult CTs, so this is a stress test rather than a clinical finding. The ensemble disagreement on these children
(E3) is 2.7 times that on the adult SZCH-X-Rays pairs (median lung-field spread 0.0076 against 0.0028) and is larger in younger children
(Spearman with age $-0.22$).

\begin{table*}[!t]
\centering
\caption{BDCXR-3257: pneumonia classification (normal, bacterial, viral), five folds over near-duplicate groups $\times$ three seeds;
pooled out-of-fold macro-AUROC, per-fold mean $\pm$ SD and the difference to the original image (group bootstrap).}
\label{tab:bdcxr_cls}
\restable[\textwidth]{bdcxr_cls}
\end{table*}

\begin{table}[!t]
\centering
\caption{BDCXR-3257: lesion audit on a fixed class-stratified subset of 300 images (1,800 sites).}
\label{tab:bdcxr_audit}
\restable{bdcxr_audit}
\end{table}

% ---------------------------------------------------------------------------------------------------------------------
\subsection{Multi-Stream Classifier Evaluation and Capacity Controls (Study~II: H6--H8)}
\label{sec:res-cls}
All 175 runs finished; Table~\ref{tab:cls_macro} gives the test metrics of the seed ensembles.

\begin{table*}[!t]
\centering
\caption{Test macro-AUROC and macro-AUPRC of the seed ensembles (five seeds; BDCXR-3257: out-of-fold over five folds). Descriptive; the
tests are in Tables~\ref{tab:cls_primary} and~\ref{tab:cls_groups}. orig: radiograph; soft: soft-tissue image; early: three images as
channels; hybrid: three streams; ctrl: radiograph in three streams.}
\label{tab:cls_macro}
\restable[\textwidth]{v3_macro}
\end{table*}

\subsubsection{Overall performance (H6, H7)}
The hybrid did not improve on the radiograph alone on any dataset (Table~\ref{tab:cls_primary}, Fig.~\ref{fig:cls_forest}). On NIH the
macro-AUROC difference was $+0.0001$ $[-0.0022,0.0024]$ (Holm $p=0.967$), so a gain larger than 0.0024 is not compatible with the data; on
VinBigData it was $+0.0024$ $[-0.0004,0.0053]$ (Holm $p=0.28$) and on BDCXR-3257 $-0.0042$ $[-0.0108,0.0021]$ (Holm $p=0.39$). The AUPRC
differences agreed in direction and none excluded zero. The hybrid did not differ from the capacity control on any dataset (Holm $p=1$).
Against channel stacking it was nominally worse on NIH ($-0.0031$ $[-0.0055,-0.0007]$, Holm $p=0.066$; AUPRC $-0.0055$
$[-0.0100,-0.0012]$) and no different elsewhere.

\begin{table*}[!t]
\centering
\caption{H6 (primary) and H7: differences of the seed ensembles with 95\,\% cluster-bootstrap intervals (\num{10000} resamples). Holm: H6
over the three datasets, H7 over its six tests.}
\label{tab:cls_primary}
\restable[\textwidth]{v3_primary}
\end{table*}

\begin{figure*}[!t]
\centering
\safefig[0.62\textwidth]{figures/fig_hypothesis_forest_densenet121.pdf}
\caption{Pre-specified contrasts of Study~II with 95\,\% cluster-bootstrap intervals. (a) H6 and H7: macro-AUROC differences per dataset.
(b) H8: differences in the mean AUROC of each finding group's classes. Filled markers: Holm $p<0.05$; right: Holm-adjusted $p$.}
\label{fig:cls_forest}
\end{figure*}

\subsubsection{Finding groups and classes (H8)}
The effect of the inputs depended on the finding group (Table~\ref{tab:cls_groups}, Fig.~\ref{fig:cls_forest}b). On NIH the soft-tissue
image raised the mean AUROC of the focal findings by $+0.0089$ $[0.0040,0.0140]$ (Holm $p=0.0024$) and the hybrid by $+0.0063$
$[0.0005,0.0117]$ (Holm $p=0.101$). Diffuse findings lost with both inputs (hybrid $-0.0052$ $[-0.0077,-0.0027]$, Holm $p=0.0024$; soft
$-0.0032$ $[-0.0057,-0.0008]$, Holm $p=0.039$), and boundary findings did not change. On VinBigData the focal estimates had the same
direction and size (hybrid $+0.0085$ $[-0.0023,0.0196]$; soft $+0.0083$ $[-0.0044,0.0202]$) but wide intervals, the diffuse estimates were
close to zero (hybrid $+0.0011$, soft $+0.0003$), and no group difference survived the correction.

\begin{table}[!t]
\centering
\caption{H8: difference in the mean per-class AUROC of each finding group, with 95\,\% cluster-bootstrap intervals; Holm within dataset
over groups and contrasts.}
\label{tab:cls_groups}
\restable{v3_groups}
\end{table}

Per class (Table~\ref{tab:cls_per_class}, Fig.~\ref{fig:cls_per_class}), two NIH classes differed between the hybrid and the radiograph
after correction: nodule gained ($+0.0129$, Holm $p=0.00121$) and infiltration lost ($-0.0160$, Holm $p<0.001$; also $-0.0128$ against
channel stacking). Mass did not change ($-0.0004$). No class differed on VinBigData or BDCXR-3257.

\begin{table*}[!t]
\centering
\caption{H8 per class: test AUROC of the radiograph and hybrid seed ensembles and the paired DeLong contrasts of the hybrid against the
radiograph and against channel stacking; Holm within dataset over classes and both contrasts.}
\label{tab:cls_per_class}
\restable[\textwidth]{v3_per_class}
\end{table*}

\begin{figure*}[!t]
\centering
\safefig[0.7\textwidth]{figures/fig_auroc_per_class_densenet121.pdf}
\caption{NIH ChestX-ray14: test AUROC per class and per finding group for each configuration, and each configuration minus the
radiograph, with 95\,\% patient-cluster bootstrap bands from 1,000 resamples (display only; the bootstrap tests of
Tables~\ref{tab:cls_primary} and~\ref{tab:cls_groups} use \num{10000}).}
\label{fig:cls_per_class}
\end{figure*}

% ---------------------------------------------------------------------------------------------------------------------
\subsection{Stream Attribution Routing and Classification Uncertainty (Study~II: H10)}
\label{sec:res-attr}

\subsubsection{Where the streams point (H10)}
Across 984 targets (164 focal) on 880 test images, Grad-CAM++ maps of the radiograph stream hit focal targets less often than other targets
(0.213 against 0.374), whereas soft-tissue stream maps hit both about equally (0.385 against 0.401) (Table~\ref{tab:cls_pointing},
Fig.~\ref{fig:cls_pointing}). The soft-minus-radiograph difference was therefore larger for focal targets ($+0.172$) than for the others
($+0.027$): interaction $+0.145$ $[0.092,0.199]$, $p=0.0002$, so H10 passes. Integrated Gradients ranked the two streams the same way
(focal: 0.580 soft-tissue against 0.505 radiograph; other: 0.621 against 0.645). Bone-stream maps rarely hit any target (0.049 and 0.167).
Fig.~\ref{fig:cls_maps} shows examples.

\begin{table*}[!t]
\centering
\caption{H10: pointing-game hit rates on the NIH expert boxes, means over targets of the per-target means over the five hybrid models;
tolerance 15\,px at 256\,px. Last row: soft-tissue minus radiograph stream (Grad-CAM++) and the interaction with its 95\,\%
image-cluster bootstrap interval.}
\label{tab:cls_pointing}
\restable[\textwidth]{v3_pointing}
\end{table*}

\begin{figure*}[!t]
\centering
\safefig[\textwidth]{figures/fig_pointing_game_densenet121.pdf}
\caption{Pointing game per NIH box class and per group. (a) Hit rate of each map. (b) Decoder routing: share of each class query's
attention on each stream. (c) Share of the Integrated-Gradients attribution per input. Brackets: number of targets.}
\label{fig:cls_pointing}
\end{figure*}

\begin{figure*}[!t]
\centering
\safefig[0.9\textwidth]{figures/fig_xai_comparison_densenet121.pdf}
\caption{Attribution maps of one hybrid model (seed 0) on NIH ChestX-ray14 test images: the three inputs, Grad-CAM++ per stream and the
decoder attention, with the expert boxes. Rows: targets (row label: the box class).}
\label{fig:cls_maps}
\end{figure*}

\subsubsection{Classification uncertainty}
Conformal risk control held the NIH false-negative rate at 0.060--0.069 on the test split for nominal 0.10, but the label sets held
6.33--6.78 of 14 labels on average (Supplementary Table~S10). On BDCXR-3257 the Mondrian sets covered 0.897--0.900 of the images at nominal
0.90 with 1.76--1.81 of 3 classes per set. In selective classification (Supplementary Table~S11 and Fig.~S13), adding the suppressor's
spread to the confidence score raised the AURC in every configuration: NIH 0.0352--0.0358 with the entropy alone against 0.0372--0.0375 with
the spread, both on the \num{25580} test images with a lung-field spread (16 of \num{25596} have an empty lung mask), and BDCXR-3257
0.1754--0.1938 against 0.2138--0.2260. The supplementary material shows every single run next to the seed ensembles (Supplementary
Fig.~S1), the conformal sets per class (Supplementary Fig.~S2) and the Integrated-Gradients maps (Supplementary Fig.~S3).

% ---------------------------------------------------------------------------------------------------------------------
\subsection{Cross-Modality Distillation, Localization and Hallucination Audits (Study~III: H11--H15)}
\label{sec:res-xm}
All 32 training runs of Study~III finished: seven CT teachers, ten projection models and 15 NIH classifiers. One zero-shot evaluation first
failed because it started before its model had been written and was repeated unchanged.

\subsubsection{CT teacher and projection models}
The seed ensemble of the CT teacher reached a test macro-AUROC of 0.821 on the 297 CT subset volumes (Table~\ref{tab:xm_transfer}): high
for cardiomegaly (0.951) and pleural effusion (0.942), low for lung lesions (0.608). The recipe was unstable across seeds: one of seven
teachers collapsed, another learned only in its last epochs, and the five teachers used ranged from 0.737 to 0.822 (Supplementary Table~S1
and Fig.~S4). On the projections, distillation lowered accuracy slightly: 0.806 for the $P_1$ ensemble against 0.814 for $P_0$, with $P_1$
below $P_0$ in every seed (Supplementary Table~S2 and Fig.~S6). Without training on real radiographs, the projection models transferred
moderately to VinBigData ($P_0$ 0.739, $P_1$ 0.732) and weakly to NIH (0.648 and 0.640).

\begin{table*}[!t]
\centering
\caption{CT teacher, projection models and zero-shot transfer: test AUROC of the seed ensembles for the six findings and their mean
(descriptive). NIH lung opacity is the approximate mapping to Infiltration.}
\label{tab:xm_transfer}
\restable[\textwidth]{v4_transfer}
\end{table*}

\subsubsection{Distillation and real radiographs (H11, H12)}
Distillation from CT did not improve the classification of real radiographs (Table~\ref{tab:xm_hyp}). On the five shared findings
$F_{\mathrm{p1}}$ differed from $F_{\mathrm{in}}$ by $-0.0001$ $[-0.0016, 0.0014]$ (Holm $p=0.912$), which excludes gains larger than
0.0014, so H11 fails. $F_{\mathrm{p1}}$ did not differ from $F_{\mathrm{p0}}$ ($-0.0008$ $[-0.0022, 0.0006]$, Holm $p=0.742$), nor
$F_{\mathrm{p0}}$ from $F_{\mathrm{in}}$ ($+0.0007$ $[-0.0008, 0.0023]$, Holm $p=0.742$), and no per-finding difference survived correction
(Supplementary Table~S3; forest plot in Supplementary Fig.~S5). Over the 14 diseases the three classifiers reached macro-AUROCs of 0.819
($F_{\mathrm{in}}$), 0.822 ($F_{\mathrm{p0}}$) and 0.821 ($F_{\mathrm{p1}}$).

\begin{table*}[!t]
\centering
\caption{Study~III hypotheses: differences of the seed ensembles with 95\,\% bootstrap intervals (\num{10000} resamples). Holm: H11 with the
two H12 contrasts. The $P_0$ row of H14 is descriptive.}
\label{tab:xm_hyp}
\restable[\textwidth]{v4_hypotheses}
\end{table*}

\subsubsection{Swin against DenseNet-121 (H13)}
DenseNet-121 exceeded $F_{\mathrm{in}}$ over the 14 diseases by 0.0049 $[0.0015, 0.0085]$ ($p=0.005$). The per-class tests place almost all
of this difference in hernia, the rarest NIH class (86 positive test images), where DenseNet-121 reached 0.901 against 0.833 (Holm
$p=0.0008$; Supplementary Table~S6 and Fig.~S12); this class alone accounts for $-0.0049$ of the 14-class mean. Infiltration also favoured
DenseNet-121 ($-0.0062$, Holm $p=0.003$), and the other 12 classes differed by less than 0.007 in either direction. H13 is not a clean
comparison of encoders: the DenseNet-121 reference was trained for twice as many epochs (20 against 10) at four times the encoder learning
rate, which may matter most for a rare class such as hernia.

\subsubsection{Cross-modal agreement (H14)}
The CT and X-ray Grad-CAM++ maps did not agree. Over 419 maps from 227 test scans, the matched CT and $P_1$ maps correlated at $-0.052$ on
average and the mismatched maps at $-0.060$, a difference of $+0.0076$ $[-0.0068, 0.0223]$ ($p=0.301$); with $P_0$ the difference was
$-0.0087$ $[-0.0228, 0.0056]$. All per-finding mean correlations were negative, and the sign of the matched-minus-mismatched difference
changed between findings (Supplementary Table~S4 and Fig.~S7). The mean maps explain the result (Fig.~\ref{fig:xm_maps}). The projected CT
maps vary almost only in the craniocaudal direction, in horizontal bands that follow the importance of individual slices, whereas the X-ray
maps of most findings concentrate along the edges of the lung box; only the cardiomegaly maps peak near the heart. Maps with such different
structure are unlikely to correlate positively in any patient.

\begin{figure*}[!t]
\centering
\safefig[\textwidth]{figures/fig_v4_maps.pdf}
\caption{Mean heat maps over the CT subset test scans positive for each finding: the teacher's Grad-CAM++ summed along the anteroposterior
axis (top) and the Grad-CAM++ maps of $P_1$ and $P_0$ on the paired projections, cropped to the lung box at $32\times32$ and scaled to sum
to one; each panel has its own colour scale. Only averages over scans are shown.}
\label{fig:xm_maps}
\end{figure*}

\subsubsection{Hallucination tests (H15)}
Rib emphasis flipped 13.7\,\% of the finding-free test images for $F_{\mathrm{in}}$ and 12.9\,\% for $F_{\mathrm{p1}}$, a difference of
$-0.0080$ $[-0.0159, -0.0003]$ ($p=0.044$; H15 is secondary and was not corrected for multiplicity). The effect is not specific to
distillation: $F_{\mathrm{p0}}$, pretrained on the same projections without the teacher, flipped 12.7\,\% and DenseNet-121 12.8\,\%
(Fig.~\ref{fig:xm_halluc}, Supplementary Table~S8). Noise separated the encoders most clearly: it flipped 78.1\,\% of the images for
DenseNet-121 and 27.3--29.3\,\% for the Swin classifiers, whereas a gamma change flipped about 6\,\% for all four. Inserted nodules were
rarely detected (sensitivity 2.8--3.2\,\%), and the score change grew with nodule size, from practically zero below 6.2\,mm to between
$+0.015$ and $+0.019$ above 8.5\,mm (Supplementary Table~S7). Between 62\,\% (DenseNet-121) and 73\,\% ($F_{\mathrm{p0}}$) of the
Grad-CAM++ mass of positive lung-finding predictions lay outside the lung fields.

\begin{figure*}[!t]
\centering
\safefig[\textwidth]{figures/fig_v4_halluc.pdf}
\caption{Hallucination tests. (a) Share of finding-free NIH test images in which some class crosses its threshold after a perturbation that
adds no disease. (b) Mean change of the Mass/Nodule score after inserting a nodule, by size tertile. (c) Share of Grad-CAM++ mass outside
the lung fields for positive lung-finding predictions.}
\label{fig:xm_halluc}
\end{figure*}

\subsubsection{Pointing game and uncertainty}
Grad-CAM++ maps of the Swin classifiers rarely pointed into the expert boxes (0.151--0.185 of the 984 targets) compared with DenseNet-121
(0.385), whereas decoder attention (0.479--0.572) and Integrated Gradients (0.615--0.645) did much better and differed little between models
(Table~\ref{tab:xm_pointing}, Supplementary Fig.~S8). All maps hit masses and nodules less often than the other boxes, and Swin Grad-CAM++
often spread along the image borders (Fig.~\ref{fig:xm_xai}). Conformal risk control kept the test false-negative rate of the three Swin
classifiers at 0.068--0.070 against a nominal 0.10, with 6.33--6.45 of the 14 labels per image (Supplementary Table~S5 and Fig.~S9).

\begin{table}[!t]
\centering
\caption{Pointing game on the 984 NIH expert boxes (164 focal: Mass, Nodule): hit rate averaged over the five seeds, with a 95\,\%
image-cluster bootstrap interval (descriptive).}
\label{tab:xm_pointing}
\restable{v4_pointing}
\end{table}

\begin{figure*}[!t]
\centering
\safefig[0.9\textwidth]{figures/fig_v4_xai_examples.pdf}
\caption{Heat maps of the seed-0 models on NIH test images with expert boxes (white): the first image of each box class among the 24
targets whose maps were saved (these contain no mass or nodule box). The NIH images are public.}
\label{fig:xm_xai}
\end{figure*}
